\documentclass[french]{article}
\usepackage{verbatim}
\usepackage{amsmath,amsfonts,amssymb}
\usepackage{pgfplots}
\usepackage[utf8]{inputenc}
\usepackage[french]{babel}
\usepackage{pgf,tikz,tkz-tab}
\usepackage{mathrsfs}
\usetikzlibrary{arrows}
\usepackage{pstricks}
\pgfplotsset{compat=newest}
\usepackage{geometry}
\geometry{hmargin=1cm,vmargin=2cm}


\begin{document}
\twocolumn
\begin{center}
\underline{\textbf{Statistiques}}
\end{center}
\section{Série statistiques et moyenne}
\section*{Exercice 1}

Calculer la moyenne de la série suivante.

   \begin{tabular}{| l | c | c | c | c | c | c |}
     \hline
     Valeurs $x_i$ & 17 & 20 & 25 & 35 & 48 & 60 \\ \hline
     Effectifs $n_i$ & 5 & 10 & 5 & 8 & 12 & 13 \\ 
     \hline
   \end{tabular}
\section*{Exercice 2}

Calculer la moyenne de la série suivante.

   \begin{tabular}{| l | c | c | c | c | c |}
     \hline
     Valeurs $x_i$ & 8 & 9 & 12 & 17 & 20 \\ \hline
     Fréquences $f_i$ & 0,13 & 0,26 & 0,08 & 0,16 & 0,37 \\ 
     \hline
    \end{tabular}
    
\section*{Exercice 3}
L'entraîneur d'un club de football a publié les statistiques des matchs de la saison. Sans calculatrice, calculer la moyenne de buts marqués par le club.

   \begin{tabular}{| l | c | c | c | c |}
     \hline
     Nombre de buts par match & 0 & 1 & 2 & 3 \\ \hline
     Fréquences & 0,35 & 0,15 & 0,4 & 0,1\\ 
     \hline
    \end{tabular}

\section*{Exercice 4}

 \begin{tabular}{| l | c | c | c | c |}
     \hline
     Valeurs $x_i$ & $x_1$ & $x_2$ & $x_2$ & $x_4$ \\ \hline
     Effectifs $n_i$ & $n_1$ & $n_2$ & $n_3$ & $n_4$ \\ 
     \hline
    \end{tabular}

\begin{enumerate}
\item On considère la série ci-dessus. Donner l'expression de la moyenne $m$ de cette série
\item Si l'on multiplie par un réel $k$ toutes les valeur quelle est la série étudiée ? Présenter les résultats sous la forme de tableau.
\item Exprimer la moyenne $m'$ de cette série.
\item En factorisant par $k$, justifier que $m'=k \times m$. 
\end{enumerate}

\section*{Exercice 5}

Dans une très petite entreprise les salaires mensuels des employés, en euros, sont les suivants: 1173 ; 1173 ; 1190 ; 1205 ; 1235 ; 3200.

\begin{enumerate}
\item Calculer la moyenne des salaires dans l'entreprise
\item Le salaire le plus élevé passe à 4850 euros. Que devient la moyenne des salaires.
\item Expliquer l'affirmation: " La moyenne est sensible aux valeurs extrêmes." 
\end{enumerate}

\section{Variance et écart-type}

\section*{Exercice 6}
Calculer la variance et l'écart-type de la série suivante. 57; 59; 62; 87; 98 et 100.
\section*{Exercice 7}
Calculer la variance et l'écart-type de la série suivante.

\begin{tabular}{| l | c | c | c | c |}
     \hline
     Valeurs $x_i$ & 10 & 12 & 17 & 20 \\ \hline
     Effectifs $n_i$ & 1 & 4 & 3 & 2  \\ 
     \hline
   \end{tabular}
\section*{Exercice 8}
On a relevé les tailles, en centimètres, des membres d'un groupe de danse folklorique.
   
   \begin{tabular}{| l | c | c | c | c | c | c |}
     \hline
     Taille (en cm) & 148 & 152 & 155 & 160 & 162 & 164 \\ \hline
     Effectifs & 1 & 3 & 5 & 2 & 5 & 3  \\ 
     \hline
   \end{tabular}
   
   \begin{enumerate}
   \item Calculer la moyenne des tailles des membres du groupe.
   \item Calculer l'écart-type.
   \end{enumerate}
   
   \section{Médiane et quartiles}
   \section*{Exercice 9}
   Déterminer la médiane et les quartiles des séries suivantes.
   \begin{enumerate}
   \item 140 ; 140 ; 148 ; 149 ; 150 ; 152 ; 152 ; 160 ; 170 ; 170 ; 175 ; 180
   \item 14 ; 15 ; 15 ; 16 ; 17 ; 17 ; 17 ; 18 ; 20 ; 20 ; 20 ; 20 ; 25
   \item \begin{tabular}{| l | c | c | c | c | c | c |}
     \hline
     Valeurs $x_i$ & 6 & 7 & 9 & 11 & 12 & 16 \\ \hline
     Effectifs $n_i$ & 2 & 3 & 4 & 8 & 5 & 3  \\ 
     \hline
   \end{tabular}
   \end{enumerate}
   
   \section*{Exercice 10}
   \begin{enumerate}
   \item Construire une série de neuf valeurs telle que $Q_1=3$, $Me=8$ et $Q_3=10$.
   \item Construire une série de cinq valeurs de moyenne 40, de médiane 20 et d'étendu 100 
   \end{enumerate}
   
   \section*{Exercice 11}
   On se donne la série suivante : 18 ; 20 ; 20 ; 22 ; 24 ; 24 ; 28 ; 30 ; 30 ; 30.
   Que deviennent la médiane et les quartiles de cette séries: 
   \begin{enumerate}
   \item si l'on ajoute 10 à toutes les valeurs ?
   \item si l'on multipllie toutes les valeurs par 10 ?
   \end{enumerate}
   
   \section*{Exercice 12}
   On donne les températures moyennes mensuelles, en dégrés Celsius, de deux villes.
  
   \begin{tabular}{| l | c | c | c | c | c | c |}
     \hline
      	& J & F & M & A & M & J \\ \hline
     Mexico & 12,4 & 14,1 & 16,2 & 17,4 & 18,4 & 17,7 \\ \hline
     Barcelone & 9,5 & 10,3 & 12,4 & 14,6 & 17,7 & 21,5  \\ 
     \hline 
   \end{tabular}\\
   
   \begin{tabular}{| l | c | c | c | c | c | c |}
     \hline
      	& J & A & S & O & N & D \\ \hline
     Mexico & 16,7 & 16,8 & 16,3 & 15,1 & 13,9 & 12 \\ \hline
     Barcelone & 24,5 & 24,3 & 21,8 & 17,6 & 13,5 & 10,3  \\ 
     \hline
   \end{tabular}\\
   
     Pour la ville de Mexico, calculer :
  
  \begin{enumerate}

  \item l'étendu des températures
  \item la température annuelle moyenne
  \item la médiane et les quartiles $Q_1$ et $Q_3$
\item[•]   Pour la ville de Barcelone, la température annuelle moyenne est de 16,48 degrés Celsius, l'étendu est de 14,8 degrés Celsius et la médiane est de 16,1 degrés Celsius, $Q_1=10,3$ et $Q_3=21,5$. Quel(s) indicateur(s) statistiques(s) permet(tent) d'affirmer que :
  \item il fait plus chaud à Barcelone qu'à Mexico ?
  \item les écarts de température sont moindres à Mexico ?
  \item dans ces deux villes, la température est supérieure à 16 degrés Celsius pendant au moins la moitié de l'année ?
  \item à Mexico, la moitié de l'année il fait approcimativement entre 14 et 17 degré Celsius.
  
  \end{enumerate}
\end{document}
